\section{AMC 10 10B 2008}
        \begin{enumerate}
        	\item (\textit{AMC 10 10B 2008}) A basketball player made $5$ baskets during a game. Each basket was worth either $2$ or $3$ points. How many different numbers could represent the total points scored by the player?


 $\mathrm{(A)}\ 2\qquad\mathrm{(B)}\ 3\qquad\mathrm{(C)}\ 4\qquad\mathrm{(D)}\ 5\qquad\mathrm{(E)}\ 6$


 
 \textbf{Answer: }E



\textbf{Solution}

The number of points could have been 10, 11, 12, 13, 14, or 15. This is because the minimum is 2*5=10 and the maximum is 3*5=15. The numbers between 10 and 15 are possible as well. Thus, the answer is $\boxed{\mathrm{(E)}\ 6}$.

 


	\item (\textit{AMC 10 10B 2008}) A $4\times 4$ block of calendar dates is shown. The order of the numbers in the second row is to be reversed. Then the order of the numbers in the fourth row is to be reversed. Finally, the numbers on each diagonal are to be added. What will be the positive difference between the two diagonal sums?


 $\begin{tabular}[t]{|c|c|c|c|} \multicolumn{4}{c}{}\\\hline 1&2&3&4\\\hline 8&9&10&11\\\hline 15&16&17&18\\\hline 22&23&24&25\\\hline \end{tabular}$


 $\textbf{(A)}\ 2 \qquad \textbf{(B)}\ 4 \qquad \textbf{(C)}\ 6 \qquad \textbf{(D)}\ 8 \qquad \textbf{(E)}\ 10$


 
 \textbf{Answer: }B



\textbf{Solution 1}

After reversing the numbers on the second and fourth rows, the block will look like this: 

 $\begin{tabular}[t]{|c|c|c|c|} \multicolumn{4}{c}{}\\\hline 1&2&3&4\\\hline 11&10&9&8\\\hline 15&16&17&18\\\hline 25&24&23&22\\\hline \end{tabular}$

 The positive difference between the two diagonal sums is then $(4+9+16+25)-(1+10+17+22)=3-1-1+3=\boxed{\textbf{(B) } 4}$.

 



\textbf{Solution 2}

Notice that at baseline the diagonals sum to the same number ($52$).  Therefore we need only compute the effect of the swap.  The positive difference between $9$ and $10$ is $1$ and the positive difference between $22$ and $25$ is $3$. Adding gives $1+3=\boxed{\textbf{(B) } 4}$

 


	\item (\textit{AMC 10 10B 2008}) Assume that $x$ is a \href{https://artofproblemsolving.com/wiki/index.php?title=Positive}{positive} \href{https://artofproblemsolving.com/wiki/index.php?title=Real\_number}{real number}. Which is equivalent to $\sqrt[3]{x\sqrt{x}}$?


 $\mathrm{(A)}\ x^{1/6}\qquad\mathrm{(B)}\ x^{1/4}\qquad\mathrm{(C)}\ x^{3/8}\qquad\mathrm{(D)}\ x^{1/2}\qquad\mathrm{(E)}\ x$


 
 \textbf{Answer: }D



\textbf{Solution 1}

$\sqrt[3]{x\sqrt{x}}=\sqrt[3]{\sqrt{x^3}}=\sqrt[6]{x^3}=x^{3/6}=x^{1/2}\ \boxed{(D)}$

 



\textbf{Solution 2}

Let x = 64. Using substitution, $\sqrt[3]{64\sqrt{64}} = \sqrt[3]{64 \cdot 8} = \sqrt[3]{512} = 8 = \sqrt{64} = 64^\frac{1}{2} = x^\frac{1}{2}$, so the answer is $\boxed{\textbf{(D)}\ x^\frac{1}{2}}$

 


	\item (\textit{AMC 10 10B 2008}) A semipro baseball league has teams with $21$ players each. League rules state that a player must be paid at least $\textdollar 15,000$ and that the total of all players' salaries for each team cannot exceed $\textdollar 700,000.$ What is the maximum possible salary, in dollars, for a single player?


 $\mathrm{(A)}\ 270,000\qquad\mathrm{(B)}\ 385,000\qquad\mathrm{(C)}\ 400,000\qquad\mathrm{(D)}\ 430,000\qquad\mathrm{(E)}\ 700,000$


 
 \textbf{Answer: }C



\textbf{Solution}

We want to find the maximum any player could make, so assume that everyone else makes the minimum possible and that the combined salaries total the maximum of $700,000$

 $700,000 = 20 * 15,000 + x$

 $x = 400,000$

 The maximum any player could make is $400,000$ dollars $\Rightarrow C$

 


	\item (\textit{AMC 10 10B 2008}) For \href{https://artofproblemsolving.com/wiki/index.php?title=Real\_number}{real numbers} $a$ and $b$, define $a \textdollar b$ $=(a-b)^2$. What is $(x-y)^2\textdollar(y-x)^2$?


 $\mathrm{(A)}\ 0\qquad\mathrm{(B)}\ x^2+y^2\qquad\mathrm{(C)}\ 2x^2\qquad\mathrm{(D)}\ 2y^2\qquad\mathrm{(E)}\ 4xy$


 
 \textbf{Answer: }A



\textbf{Solution}

Since $(-a)^2 = a^2$, it follows that $(x-y)^2 = (y-x)^2$, and 
 \[(x-y)^2 * (y-x)^2 = [(x-y)^2 - (y-x)^2]^2 = [(x-y)^2 - (x-y)^2]^2 = 0\ \mathrm{(A)}.\]

 


	\item (\textit{AMC 10 10B 2008}) Points $B$ and $C$ lie on $AD$. The length of $AB$ is $4$ times the length of $BD$, and the length of $AC$ is $9$ times the length of $CD$. The length of $BC$ is what fraction of the length of $AD$?


 $\mathrm{(A)}\ 1/36\qquad\mathrm{(B)}\ 1/13\qquad\mathrm{(C)}\ 1/10\qquad\mathrm{(D)}\ 5/36\qquad\mathrm{(E)}\ 1/5$


 
 \textbf{Answer: }C



\textbf{Solution 1}

Let $CD = 1$. Then $AB = 4(BC + 1)$ and $AB + BC = 9\cdot1$. From this system of equations, we obtain $BC = 1$. Adding $CD$ to both sides of the second equation, we obtain $AD = AB + BC + CD = 9 + 1 = 10$.  Thus, $\frac{BC}{AD} = \frac{1}{10} \implies\text{(C)}$

 



\textbf{Solution 2}

Let $x = BD$ and $y = CD$. Therefore, $AB = 4x$ and $AC = 9y$, as shown in the diagram(the labels on the bottom are for that line segment while the labels on the top are from one point to the left to one point to the right).  

 
                    \begin{figure}[H]
                        \centering
                        \includegraphics[width=0.5\textwidth]{Solution/AMC_10_10B_2008_6_Solution_Figure_1.png}
                    \end{figure}
                    

 From this, we can see that $AD = 10y = 5x$, and since $BC = BD - CD = x-y$. Now, our ratio is $\frac{x-y}{AD}$. We can split this into 2 fractions: $\frac{x}{AD} - \frac{y}{AD} = \frac{x}{5x} - \frac{y}{10y} = \frac{1}{5} - \frac{1}{10} =  \boxed{\textbf{(C)}\ \frac{1}{10}}$

 


	\item (\textit{AMC 10 10B 2008}) An equilateral triangle of side length $10$ is completely filled in by non-overlapping equilateral triangles of side length $1$. How many small triangles are required?


 $\mathrm{(A)}\ 10\qquad\mathrm{(B)}\ 25\qquad\mathrm{(C)}\ 100\qquad\mathrm{(D)}\ 250\qquad\mathrm{(E)}\ 1000$


 
 \textbf{Answer: }C



\textbf{Solution 1}

The area of the large triangle is $\frac{10^2\sqrt3}{4}$, while the area of each small triangle is $\frac{1^2\sqrt3}{4}$. Dividing these two quantities results in $100$, therefore $\boxed{100} \mathrm{(C)}$ small triangles can fill the large one without overlap.

 



\textbf{Solution 2-4}


                    \begin{figure}[H]
                        \centering
                        \includegraphics[width=0.5\textwidth]{Solution/AMC_10_10B_2008_7_Solution_Figure_1.png}
                    \end{figure}
                    

 The number of triangles is $1+3+\dots+19 = \boxed{100}$.

 Also, another way to do it is to notice that as you go row by row (from the bottom), the number of triangles decrease by 2 from 19, so we have:$19+17+15...+3+1 = \frac{19+1}{2}\cdot 10 = \boxed{100}$

 A fourth solution is to notice that the small triangles are similar to the large triangle as they are both equilateral. Therefore, the ratio of their areas is the square of the ratio of their side lengths.  Hence the ratio of their areas is $(1/10)^2=1/100$, so the answer is $\boxed{100}$.

 



\textbf{Solution 5}

The side length of the large equilateral triangle is $10$. The number of equilateral triangles with side length $1$ that can fill it is $10^2$. This is equal to the sum of the first $10$ odd integers. The sum of the first $n$ odd integers is $n^2$. In general, an equilateral triangle with side length $n$ can be filled by $n^2$ triangles. The answer is $100\ \mathrm{(C)}$.

 


	\item (\textit{AMC 10 10B 2008}) A class collects $\textdollar50$ to buy flowers for a classmate who is in the hospital. Roses cost $\textdollar3$ each, and carnations cost $\textdollar2$ each. No other flowers are to be used. How many different bouquets could be purchased for exactly $\textdollar50$?


 $\mathrm{(A)}\ 1 \qquad \mathrm{(B)}\ 7 \qquad \mathrm{(C)}\ 9 \qquad \mathrm{(D)}\ 16 \qquad \mathrm{(E)}\ 17$


 
 \textbf{Answer: }C



\textbf{Solution 1}

The cost of a rose is odd, hence we need an even number of roses. Let there be $2r$ roses for some $r\geq 0$. Then we have $50-3\cdot 2r = 50-6r$ dollars left. We can always reach the sum exactly $50$ by buying $(50-6r)/2 = 25-3r$ carnations. Of course, the number of roses must be such that the number of carnations is non-negative. We get the inequality $25-3r \geq 0$ which solves to $r\leq 8 \frac13$. $r$ must be an integer, so there are $\boxed{9 \text{ (C)}}$ possible values of $r$, and each gives us one solution.

 



\textbf{Solution 2}

Let $x$ and $y$ be the number of roses and carnations bought. The equation should be $3x+2y = 50$. Since $50$ is an even number, the product of $3x$ must be even and smaller than $50$. You can try nonnegative even integers for $x$ and you will end up with the numbers $0$, $2$, $4$, $6$, $8$, $10$, $12$, $14$, and $16$. There are $9$ numbers in total, so the answer is  $\boxed{9 \text{ (C)}}$.

 



\textbf{Solution 3}

Let $r$ represent the number of roses, and let $c$ represent the number of carnations. Then, we get the linear Diophantine equation,$3r+2c=50$.Using the Euclidean algorithm, we get the initial solutions to be $r_0=50$ and $c_0=-50$, meaning the complete solution will be,$r=50+\frac{2}{\gcd(2,3)}$ $k=50+2k$, $c=-50-\frac{3}{\gcd(2,3)}k=-50-3k$

 The solution range for which both $r$ and $c$ are positive is $17$ $\leq k$ $\leq$ $25$. There are $\boxed{9 \text{ (C)}}$ possible values for $k$.

 


	\item (\textit{AMC 10 10B 2008}) A quadratic equation $ax^2 - 2ax + b = 0$ has two real solutions. What is the average of these two solutions?


 $\mathrm{(A)}\ 1\qquad\mathrm{(B)}\ 2\qquad\mathrm{(C)}\ \frac ba\qquad\mathrm{(D)}\ \frac{2b}a\qquad\mathrm{(E)}\ \sqrt{2b-a}$


 
 \textbf{Answer: }A



\textbf{Solution 1}

Dividing both sides by $a$, we get $x^2 - 2x + b/a = 0$. By Vieta's formulas, the sum of the roots is $2$, therefore their average is $1\Rightarrow \boxed{A}$.

 



\textbf{Solution 2}

We know that for an equation $ax^2 + bx + c = 0$, the sum of the roots is $\frac{-b}{a}$. This means that the sum of the roots for $ax^2 - 2ax + b = 0$ is $\frac{2a}{a}=2$. The average is the sum of the two roots divided by two, so the average is $\frac22 = 1 \Rightarrow \boxed{A}$.

 



\textbf{Solution 3}

Using the quadratic formula(assuming $a=1$),we get $x = \frac{2 \pm \sqrt{4-4b}}{2}$. Let $y = \sqrt{4-4b}$. Simplifying, we get $1 \pm \frac{y}{2}$. That means the sum of the solutions is $(1+\frac{y}{2}) + (1-\frac{y}{2}) = 2$, so the average is $2/2 = \boxed{\textbf{(A) } 1}$.

 


	\item (\textit{AMC 10 10B 2008}) Points $A$ and $B$ are on a circle of radius $5$ and $AB=6$. Point $C$ is the \href{https://artofproblemsolving.com/wiki/index.php?title=Midpoint}{midpoint} of the minor arc $AB$. What is the length of the line segment $AC$?


 $\mathrm{(A)}\ \sqrt{10}\qquad\mathrm{(B)}\ \frac{7}{2}\qquad\mathrm{(C)}\ \sqrt{14}\qquad\mathrm{(D)}\ \sqrt{15}\qquad\mathrm{(E)}\ 4$


 
 \textbf{Answer: }A



\textbf{Solution}

Let the center of the circle be $O$, and let $D$ be the intersection of $\overline{AB}$ and $\overline{OC}$ (then $D$ is the midpoint of $\overline{AB}$). $OA=OB=5$, since they are both radii of the circle.

 By the \href{https://artofproblemsolving.com/wiki/index.php?title=Pythagorean\_Theorem}{Pythagorean Theorem}, $OD = \sqrt{OA^2 - DA^2} = 4$, and by subtraction, $CD=OC - OD = 1$.

 Using the Pythagorean Theorem again, $AC= \sqrt{AD^2 + CD^2} = \sqrt{3^2+1^2}=\sqrt{10} \Longrightarrow \textbf{(A)}$.

 
                    \begin{figure}[H]
                        \centering
                        \includegraphics[width=0.5\textwidth]{Solution/AMC_10_10B_2008_10_Solution_Figure_1.png}
                    \end{figure}
                    

 


	\item (\textit{AMC 10 10B 2008}) Suppose that $(u_n)$ is a \href{https://artofproblemsolving.com/wiki/index.php?title=Sequence}{sequence} of real numbers satisfying $u_{n+2}=2u_{n+1}+u_n$,


 and that $u_3=9$ and $u_6=128$. What is $u_5$?


 $\mathrm{(A)}\ 40\qquad\mathrm{(B)}\ 53\qquad\mathrm{(C)}\ 68\qquad\mathrm{(D)}\ 88\qquad\mathrm{(E)}\ 104$


 
 \textbf{Answer: }B



\textbf{Solution 1}

If we plug in $n=4$, we get 

 $128=2u_5+u_4.$

 By plugging in $n=3$, we get 

 $u_5=2u_4+9.$

 This is a system of two equations with two unknowns. Multiplying the second equation by 2 and substituting into the first equation gives $128=5u_4+18 \Longrightarrow u_4=22$, therefore $u_5=\frac{128-22}{2}=53 \longrightarrow \textbf{\boxed{(B)}}$.

 



\textbf{Solution 2}

Suppose the first 2 terms are x and y, then we can list the rest of the terms by the given definition.

 $x , y, 2y+x, 5y+2x, 12y+5x, 29y+12x$

 Then from this we get the equations $29y+12x=128$ and $2y+x=9$. Solving these we get $y=4$ and $x=1$. Then we find the 5th term by substituting $y=4$ and $x=1$ into $12y+5x$ and we get $\boxed{(B)}$

 


	\item (\textit{AMC 10 10B 2008}) Postman Pete has a pedometer to count his steps. The pedometer records up to 99999 steps, then flips over to 00000 on the next step. Pete plans to determine his mileage for a year. On January 1 Pete sets the pedometer to 00000. During the year, the pedometer flips from 99999 to 00000 forty-four 
times. On December 31 the pedometer reads 50000. Pete takes 1800 steps per mile. Which of the following is closest to the number of miles Pete walked during the year?


 $\mathrm{(A)}\ 2500\qquad\mathrm{(B)}\ 3000\qquad\mathrm{(C)}\ 3500\qquad\mathrm{(D)}\ 4000\qquad\mathrm{(E)}\ 4500$


 
 \textbf{Answer: }A



\textbf{Solution}

Every time the pedometer flips from $99999$ to

 $00000$ Pete has walked $100000$ steps.

 So, if the pedometer flipped $44$ times

 Pete walked $44*100000+50000=4450000$ steps.

 Dividing by $1800$ steps per mile gives $2472.\overline{2}$

 This is closest to answer $\boxed{A}$.

 


	\item (\textit{AMC 10 10B 2008}) For each positive integer $n$, the mean of the first $n$ terms of a sequence is $n$. What is the 2008th term of the sequence?


 $\mathrm{(A)}\ {{{2008}}} \qquad \mathrm{(B)}\ {{{4015}}} \qquad \mathrm{(C)}\ {{{4016}}} \qquad \mathrm{(D)}\ {{{4,030,056}}} \qquad \mathrm{(E)}\ {{{4,032,064}}}$


 
 \textbf{Answer: }B



\textbf{Solution 1}

Since the mean of the first $n$ terms is $n$, the sum of the first $n$ terms is $n^2$. Thus, the sum of the first $2007$ terms is $2007^2$ and the sum of the first $2008$ terms is $2008^2$. Hence, the $2008^{\text{th}}$ term of the sequence is $2008^2-2007^2=(2008+2007)(2008-2007)=4015\Rightarrow \boxed{\textbf{(B) 4015}}$

 Note that $n^2$ is the sum of the first n odd integers.

 


 



\textbf{Solution 2}

Let $a_1, a_2, a_3, \cdots, a_n$ be the terms of the sequence. We know $\frac{a_1 + a_2 + a_3 + \cdots + a_n}{n} = n$, so we must have $a_1 + a_2 + a_3 + \cdots + a_n = n^2$. The sum of consecutive odd numbers down to $1$ is a perfect square, if you don't believe me, try drawing squares with the sum, so $a_1 = 1, a_2 = 3, a_3 = 5, \cdots , a_n = 2(n-1) + 1$, so the answer is $a_{2008} = 2(2007) + 1 = \boxed{\textbf{(B) 4015}}$.

 



\textbf{Solution 3}

Let the mean be $\frac{(a)+(a+d)+(a+2d)+...+(a+(n-1)) \cdot d)}{n}$$=\frac{n \cdot a}{n} + \frac{(1+2+3+...+(n-1)) \cdot d}{n}$$=a + \frac{n \cdot (n-1) \cdot d}{2n}$$=a+ \frac{(n-1) \cdot d}{2}$

 Note that this is also equal to n

 $a+ \frac{(n-1) \cdot d}{2}=n$$\therefore 2a+ (n-1) \cdot d=2n$

 1st term + nth term $=2n=2 \cdot 2008=4016$Now note that, from previous solutions, the first term is 1, hence the 2008th term is $4016-1=\boxed{\textbf{(B) 4015}}$

 



\textbf{Solution 4 (Using Answer Choices)}

From inspection, we see that the sum of the sequence is $n^2$. We also notice that $n^2$ is the sum of the first $n$ odd integers. Because $4015$ is the only odd integer, $\boxed{\textbf{(B) 4015}}$ is the answer.

 


 


	\item (\textit{AMC 10 10B 2008}) Triangle $OAB$ has $O=(0,0)$, $B=(5,0)$, and $A$ in the first quadrant. In addition, $\angle ABO=90^\circ$ and $\angle AOB=30^\circ$. Suppose that $OA$ is rotated $90^\circ$ counterclockwise about $O$. What are the coordinates of the image of $A$? 


 $\mathrm{(A)}\ \left( - \frac {10}{3}\sqrt {3},5\right) \qquad \mathrm{(B)}\ \left( - \frac {5}{3}\sqrt {3},5\right) \qquad \mathrm{(C)}\ \left(\sqrt {3},5\right) \qquad \mathrm{(D)}\ \left(\frac {5}{3}\sqrt {3},5\right) \qquad \mathrm{(E)}\ \left(\frac {10}{3}\sqrt {3},5\right)$


 
 \textbf{Answer: }B



\textbf{Solution 1}

Since $\angle ABO=90^\circ$, and $\angle AOB=30^\circ$, we know that this triangle is one of the \href{https://artofproblemsolving.com/wiki/index.php?title=Special\_Right\_Triangles}{Special Right Triangles}.

 We also know that $B$ is $(5,0)$, so $B$ lies on the x-axis. Therefore, $OB = 5$.

 Then, since we know that this is a Special Right Triangle ($30$-$60$-$90$ triangle), we can use the proportion 
 \[\frac{5}{\sqrt 3}=\frac{AB}{1}\] to find $AB$.

 We find that 
 \[AB=\frac{5\sqrt 3}{3}\]

 That means the coordinates of $A$ are $\left(5,\frac{5\sqrt 3}3\right)$.

 Rotate this triangle $90^\circ$ counterclockwise around $O$, and you will find that $A$ will end up in the second quadrant with the coordinates$\left( -\frac{5\sqrt 3}3, 5\right)$, or $\boxed{(B)}$

 Note: To better visualize this, one can sketch a diagram.

 


 


 



\textbf{Solution 2}

As $\angle ABO=90^\circ$ and $A$ is in the first quadrant, we know that the $x$ coordinate of $A$ is $5$. We now need to pick a positive $y$ coordinate for $A$ so that we'll have $\angle AOB=30^\circ$. 

 By the \href{https://artofproblemsolving.com/wiki/index.php?title=Pythagorean\_theorem}{Pythagorean theorem} we have $AO^2 = AB^2 + BO^2 = AB^2 + 25$. 

 By the definition of \href{https://artofproblemsolving.com/wiki/index.php?title=Sine}{sine}, we have $\frac{AB}{AO} = \sin AOB = \sin 30^\circ = \frac 12$, hence $AO=2\cdot AB$.

 Substituting into the previous equation, we get $AB^2 = \frac{25}3$, hence $AB=\frac{5\sqrt 3}3$.

 This means that the coordinates of $A$ are $\left(5,\frac{5\sqrt 3}3\right)$. 

 After we rotate $OA$ $90^\circ$ counterclockwise about $O$, it will be in the second quadrant and have the coordinates $\boxed{ \left( -\frac{5\sqrt 3}3, 5\right) }$.

 



\textbf{Solution 3 (cheese)}

Notice that the answer must be in the form $(-a,b)$, so $\textbf{(C)}, \textbf{(D)}$ and $\textbf{(E)}$ will not be the answer

 Using graphing paper and a protractor, the answer can be estimated to be $\boxed{\textbf{(B)} \left( - \frac {5}{3}\sqrt {3},5\right)}$

 


	\item (\textit{AMC 10 10B 2008}) How many right triangles have integer leg lengths $a$ and $b$ and a hypotenuse of length $b+1$, where $b<100$?


 $\mathrm{(A)}\ 6\qquad\mathrm{(B)}\ 7\qquad\mathrm{(C)}\ 8\qquad\mathrm{(D)}\ 9\qquad\mathrm{(E)}\ 10$


 
 \textbf{Answer: }A



\textbf{Solution}

By the Pythagorean theorem, $a^2+b^2=b^2+2b+1$

 This means that $a^2=2b+1$.

 We know that $a,b>0$ and that $b<100$.

 We also know that $a^2$ is odd and thus $a$ is odd, since the right side of the equation is odd. $2b$ is even. $2b+1$ is odd.

 So $a=1,3,5,7,9,11,13$, but if $a=1$, then $b=0$. Thus $a\neq1.$

 $a=3,5,7,9,11,13$

 The answer is $\boxed{A}$.

 


	\item (\textit{AMC 10 10B 2008}) Two fair coins are to be tossed once. For each head that results, one fair die is to be rolled. What is the probability that the sum of the die rolls is odd? (Note that if no die is rolled, the sum is $0$.)


 $\mathrm{(A)}\ {{{\frac{3} {8}}}} \qquad \mathrm{(B)}\ {{{\frac{1} {2}}}} \qquad \mathrm{(C)}\ {{{\frac{43} {72}}}} \qquad \mathrm{(D)}\ {{{\frac{5} {8}}}} \qquad \mathrm{(E)}\ {{{\frac{2} {3}}}}$


 
 \textbf{Answer: }A



\textbf{Solution 1 (Casework)}

We consider 3 cases based on the outcome of the coin:

 Case 1, 0 heads:The probability of this occurring on the coin flip is $\frac{1} {4}$. The probability that 0 rolls of a die will result in an odd sum is $0$. 

 Case 2, 1 head:The probability of this case occurring is $\frac{1} {2} \cdot \frac {1} {2} \cdot 2 = \frac {1} {2}.$ The probability that 1 die results in an odd number is $\frac{1} {2}$.

 Case 3, 2 heads:The probability of this occurring is $\frac{1} {4}$. The probability that 2 dice result in an odd sum is $\frac{1} {2}$, because regardless of what we throw on the first die, we have $\frac{1} {2}$ probability that the second die will have the opposite parity.

 Thus, the probability of having an odd sum rolled is $\frac{1} {4} \cdot 0 + \frac{1} {2} \cdot \frac{1} {2} + \frac{1} {4} \cdot \frac{1} {2}=\frac{3} {8}\Rightarrow \boxed{A}$

 



\textbf{Solution 2 (possibly slightly faster)}

We use complementary counting or subtracting $P(\text{Even})$ from $1$. We use casework now.

 Case $1$: $2$ Tails. $2$ tails occur with probability $\frac{1}{4}$, but we will always get an even number, so the overall probability to get an even sum is $\frac{1}{4}$.

 Case $2$: $1$ Tail: This event occurs with probability $\frac{1}{2}$ and the probability we get an even is $\frac{1}{2}$, so the overall probability to get an even, in this case, is also $\frac{1}{4}$.

 We know $P\text{(Even)}$ is greater than $\frac{1}{2}$, so $P\text{(Odd)}$ is less than $\frac{1}{2}$.

 Only $\boxed{\text{(A)}\frac{3}{8}}$ is less than $\frac{1}{2}$.

 



\textbf{Solution 3}

When rolling one or two dice, the probability of the sum being odd is equivalent to the probability of the sum being even. However, if we flip no heads, meaning we roll no dice, the sum is always even. Hence, the probability that the sum of the dice is even is ( 1 - the probability of flipping no heads ) / 2. Solving, we get that the answer is (1 - $\frac{1}{4}$) / 2 = $\frac{3}{8}$.

 


 


 



\textbf{Solution 4 (Very Fast)}

0 heads: If you are not rolling the dice, then it is an even number (0). 1 heads: If you are rolling 1 dice, the probability is 1/2, the same as odd. 2 heads: When you are rolling 2 dices, there are 2 choices for odd (O,E) (E,O) while there are also 2 choices for even (E,E) (O,O). In total, the chances of rolling an even is more than rolling an odd (More, Same, Same). So, the answer must be less than half and the only answer choice that satisfies this is answer choice (A).

 


 \begin{verbatim}
	~yanchudeng
\end{verbatim}




	\item (\textit{AMC 10 10B 2008}) A poll shows that $70\%$ of all voters approve of the mayor's work. On three separate occasions a pollster selects a voter at random. What is the probability that on exactly one of these three occasions the voter approves of the mayor's work?


 $\mathrm{(A)}\ {{{0.063}}} \qquad \mathrm{(B)}\ {{{0.189}}} \qquad \mathrm{(C)}\ {{{0.233}}} \qquad \mathrm{(D)}\ {{{0.333}}} \qquad \mathrm{(E)}\ {{{0.441}}}$


 
 \textbf{Answer: }B



\textbf{Solution 1}

Letting Y stand for a voter who approved of the work, and N stand for a person who didn't approve of the work, the pollster could select responses in $3$ different ways: $\text{YNN, NYN, and NNY}$. The probability of each of these is $(0.7)(0.3)^2=0.063$. Thus, the answer is $3\cdot0.063=0.189$, or $\boxed{B}$

 



\textbf{Solution 2}

In more concise terms, this problem is an extension of the binomial distribution. We find the number of ways only 1 person approves of the mayor multiplied by the probability 1 person approves and 2 people disapprove: ${3\choose 1} \cdot(0.7)^1\cdot(1-0.7)^{(3-1)}=3\cdot0.7\cdot0.09=\boxed{\mathrm{(B)}\ {{{0.189}}}}$

 



\textbf{Solution 3 (combinatorics)}

The probability of getting the first voter to approve is $\frac{7}{10} * \frac{3}{10} * \frac{3}{10}$.

 This first voter, using combinations, can be arranged in 3 choose 1 ways, which simplifies into 3 ways.

 Multiplying $3$ by $\frac{7}{10} * \frac{3}{10} * \frac{3}{10}$ gives (B).

 PS: For a more challenging practice problem akin to this one, look to 2012 AMC 12A Problem 11.

 


	\item (\textit{AMC 10 10B 2008}) Bricklayer Brenda would take nine hours to build a chimney alone, and Bricklayer Brandon would take $10$ hours to build it alone. When they work together, they talk a lot, and their combined output decreases by $10$ bricks per hour. Working together, they build the chimney in $5$ hours. How many bricks are in the chimney?


 $\mathrm{(A)}\ 500\qquad\mathrm{(B)}\ 900\qquad\mathrm{(C)}\ 950\qquad\mathrm{(D)}\ 1000\qquad\mathrm{(E)}\ 1900$


 
 \textbf{Answer: }B



\textbf{Solution}

Let $x$ be the number of bricks in the chimney. The work done is the rate multiplied by the time. 

 Using $w = rt$, we get $x = \left(\frac{x}{9} + \frac{x}{10} - 10\right)\cdot(5)$. Solving for $x$, we get $900$, or $\boxed{B}$.

 



\textbf{Solution2 (only use as last resort)}

We take a look at the answer choices and see that only B is a multiple of both 9 and 10.-mathlover1205

 


	\item (\textit{AMC 10 10B 2008}) A cylindrical tank with radius $4$ feet and height $9$ feet is lying on its side. The tank is filled with water to a depth of $2$ feet. What is the volume of water, in cubic feet?


 $\mathrm{(A)}\ 24\pi - 36 \sqrt {2} \qquad \mathrm{(B)}\ 24\pi - 24 \sqrt {3} \qquad \mathrm{(C)}\ 36\pi - 36 \sqrt {3} \qquad \mathrm{(D)}\ 36\pi - 24 \sqrt {2} \qquad \mathrm{(E)}\ 48\pi - 36 \sqrt {3}$


 
 \textbf{Answer: }E



\textbf{Solution}

Any vertical cross-section of the tank parallel with its base looks as follows:
                    \begin{figure}[H]
                        \centering
                        \includegraphics[width=0.5\textwidth]{Solution/AMC_10_10B_2008_19_Solution_Figure_1.png}
                    \end{figure}
                    

 The volume of water can be computed as the height of the tank times the area of the shaded part.

 Let $\theta$ be the size of the smaller angle $DAC$. We then have $\cos\theta = \frac{AD}{AC}=\frac 12$, hence $\theta=60^\circ$.

 The figure is symmetrical, so the angle $CAB$ has size $2\cdot 60^\circ = 120^\circ$. Hence the shaded part consists of $\frac{120^\circ}{360^\circ} = \frac 13$ of the circle, minus the area of the triangle $ABC$.

 Using the \href{https://artofproblemsolving.com/wiki/index.php?title=Pythagorean\_theorem}{Pythagorean theorem} we can compute that $CD=\sqrt{AC^2-AD^2}=\sqrt{16-4}=2\sqrt 3$. Thus $BC=4\sqrt 3$, and the area of the triangle $ABC$ is $\frac {2 \cdot 4\sqrt 3} 2 = 4\sqrt 3$.

 The area of the shaded part is then $\frac{4^2\pi}3 - 4\sqrt 3$, and the volume of water is $9\cdot\left(\frac{4^2\pi}3 - 4\sqrt 3\right) = \boxed{48\pi - 36\sqrt 3}$. The answer is $\text{E}$.

 Alternatively, we can solve for the length of $\overline{DC}$, as it is congruent to a 30-60-90 triangle by the SSS congruency theorem. ~Williamgolly

 


	\item (\textit{AMC 10 10B 2008}) The faces of a cubical die are marked with the numbers $1$, $2$, $2$, $3$, $3$, and $4$. The faces of another die are marked with the numbers $1$, $3$, $4$, $5$, $6$, and $8$. Both dice are thrown. What is the probability that the sum of the top two numbers will be $5$, $7$, or $9$?                   


 $\mathrm{(A)}\ 5/18\qquad\mathrm{(B)}\ 7/18\qquad\mathrm{(C)}\ 11/18\qquad\mathrm{(D)}\ 3/4\qquad\mathrm{(E)}\ 8/9$


 
 \textbf{Answer: }B



\textbf{Solution 1}

One approach is to write a table of all $36$ possible outcomes, do the sums, and count good outcomes.

 \begin{verbatim}
	1  3  4  5  6  8   ------------------1 |  2  4  5  6  7  92 |  3  5  6  7  8 102 |  3  5  6  7  8 103 |  4  6  7  8  9 113 |  4  6  7  8  9 114 |  5  7  8  9 10 12
\end{verbatim}

We see that out of $36$ possible outcomes, $4$ give the sum of $5$, $6$ the sum of $7$, and $4$ the sum of $9$, hence the resulting probability is $\frac{4+6+4}{36}=\frac{14}{36}=\boxed{\frac{7}{18}}$.

 



\textbf{Solution 2}

Each die is equally likely to roll odd or even, so the probability of an odd sum is $\frac{1}{2}$.The possible odd sums are $3, 5, 7, 9, 11$.

 So we can find the probability of rolling $3$ or $11$ instead and just subtract that from $\frac{1}{2}$, which seems easier. 

 Without writing out a table, we can see that there are 2 ways to make $3$, and 2 ways to make $11$, for a probability of $\frac{4}{36}$.

 $\frac{1}{2}-\frac{4}{36}=\frac{14}{36}=\boxed{\frac{7}{18}}$.

 



\textbf{Solution 3}

The outcome of the first dice can be $1$, $2$, $3$, or $4$. For each of these $4$ cases, we can find the possible outcomes for the second dice that makes the sum of the top two numbers $5$, $7$, or $9$, and then calculate the respective probabilities. 

 \textbf{Case 1 - the first dice is 1:} the outcome of the second dice can be $4$, $6$, or $8$. There is a $\frac{1}{6}$ probability of rolling a $1$ with the first dice and a $\frac{1}{2}$ probability of rolling a $4$, $6$, or $8$ with the second dice. $\frac{1}{6}*\frac{1}{2} = \frac{1}{12}.$

 \textbf{Case 2 - the first dice is 2:} the outcome of the second dice can be $3$ or $5$. There is a $\frac{1}{3}$ probability of rolling a $2$ with the first dice and a $\frac{1}{3}$ probability of rolling a $3$ or $5$ with the second dice. $\frac{1}{3}*\frac{1}{3} = \frac{1}{9}.$

 \textbf{Case 3 - the first dice is 3:} the outcome of the second dice can be $4$ or $6$. There is a $\frac{1}{3}$ probability of rolling a $3$ with the first dice and a $\frac{1}{3}$ probability of rolling a $4$ or $6$ with the second dice. $\frac{1}{3}*\frac{1}{3} = \frac{1}{9}.$

 \textbf{Case 4 - the first dice is 4:} the outcome of the second dice can be $1$, $3$, or $5$. There is a $\frac{1}{6}$ probability of rolling a $4$ with the first dice and a $\frac{1}{2}$ probability of rolling a $1$, $3$, or $5$ with the second dice. $\frac{1}{6}*\frac{1}{2} = \frac{1}{12}.$

 $\frac{1}{12} + \frac{1}{9} + \frac{1}{9} + \frac{1}{12} = \frac{14}{36} = \frac{7}{18}$, so the answer is $\boxed{\textbf{(B) } \frac{7}{18}}$. ~azc1027

 


	\item (\textit{AMC 10 10B 2008}) Ten chairs are evenly spaced around a round table  and numbered clockwise from $1$ through $10$. Five married couples are to sit in the chairs with men and women alternating, and no one is to sit either next to or across from his/her spouse. How many seating arrangements are possible?


 $\mathrm{(A)}\ 240\qquad\mathrm{(B)}\ 360\qquad\mathrm{(C)}\ 480\qquad\mathrm{(D)}\ 540\qquad\mathrm{(E)}\ 720$


 
 \textbf{Answer: }C



\textbf{Solution}

For the first man, there are $10$ possible seats. For each subsequent man, there are $4$, $3$, $2$, or $1$ possible seats. After the men are seated, there are only two possible arrangements for the five women. The answer is $10\cdot 4\cdot 3\cdot 2\cdot 1\cdot 2 = \boxed{(\text{C}) 480}$.

 



\textbf{Solution 2}

Label the seats ABCDEFGHIJ, where A is the top seat. The first man has $10$ possible seats. WLOG, assume he is in seat A in the diagram. Then, his wife can be in one of two seats, namely D or H. WLOG, assume she is in seat D. Now, in each structurally distinct solution we find, we know that there are $4! = 24$ ways to arrange the 4 other couples. Let there be x structurally distinct solutions under these conditions. We know the answer must be $10\cdot 2\cdot 24\cdot x = 480x$ possible seating arrangements, and x is a nonnegative integer. There is only one answer that is a multiple of $480$. So, our answer is $\boxed{(\text{C}) 480}$.

 


	\item (\textit{AMC 10 10B 2008}) Three red beads, two white beads, and one blue bead are placed in line in random order. What is the probability that no two neighboring beads are the same color?


 $\mathrm{(A)}\ 1/12\qquad\mathrm{(B)}\ 1/10\qquad\mathrm{(C)}\ 1/6\qquad\mathrm{(D)}\ 1/3\qquad\mathrm{(E)}\ 1/2$


 
 \textbf{Answer: }C



\textbf{Solution 1 (Order Matters)}

There are two ways to arrange the red beads, where $R$ represents a red bead and $-$ represents a blank space.

 $1. R - R - R -$

 $2. R - - R - R$

 In the first, there are three ways to place a bead in the first free space, two for the second free space, and one for the third, so there are $6$ arrangements.In the second, a white bead must be placed in the third free space, so there are two possibilities for the third space, two for the second, and one for the first. That makes $4$ arrangements. There are $6+4=10$ arrangements in total. The two cases above can be reversed, so we double $10$ to get $20$ arrangements. Also, in each case, there are three ways to place the first red bead, two for the second, and one for the third, so we multiply by $6$ to get $20\cdot 6 = 120$ arrangements. There are $6! = 720$ total arrangements so the answer is $120/720 = \boxed{1/6}$.

 



\textbf{Solution 2 (Order Doesn't Matter)}

If you list the cases for the red balls and blanks, these are four scenarios:

 $1. R - - R - R$

 $2. R - R - - R$

 $3. R - R - R -$

 $4. - R - R - R$

 In cases 1 and 2, the white balls must go in the blank surrounded on either side by the red balls. Like this:$1. R - - R W R$$2. R W R - - R$However, now with only one white ball and one blue ball left, there are two ways to order them in both cases. So for a total of $2 + 2 = 4$ possibilities.

 In the last two cases, the two white balls and the blue ball can go anywhere in those three blanks, because they are separated by a red ball. So there are $3$ ways for each case. This adds up to a total of $6$ possibilities for the last two cases.

 Adding them up, we get $10$ total orderings.

 There are $\frac{6!}{3! * 2! * 1!} = 60$ total orderings.

 So the answer is $\frac{10}{60} = \boxed{1/6}$.

 



\textbf{Solution 3}

There are 3 cases:

 \textbf{Case 1} There is a block of $RWRWR$. The $B$ can go in any slot between two letters. 6 ways.

 \textbf{Case 2} There is a $RBR$. Then, there is a $WRW$, and there are 2 orderings.

 \textbf{Case 3} There is a $WBW$. Then, there is a $RWR$, and there are 2 orderings.

 There are 60 ways total, and 10 valid ones, giving $\boxed{\frac{1}{6}}$ as the answer.

 



\textbf{Solution 4}

There are $\frac{6!}{3!\cdot2!\cdot1!}=60$ total orderings.

 Suppose we order the red and white beads first. If these two colors are ordered first, we must make sure that no neighboring beads are the same color, or only one pair of neighboring beads are the same color. There are five possible orderings:

 $R\ W R\ W R$

 $R\ R\ W R\ W$

 $W R\ R\ W R$

 $R\ W R\ R\ W$

 $W R\ W R\ R$

 For the first case, there are $6$ possible places we can put the blue bead. For the other 4 cases, there is only one place we can put the blue bead, which is between the pair of red beads. The desired probability is $\frac{6+4(1)}{60}=\frac{10}{60}=\boxed{\mathrm{(C) }\frac16}$

 


	\item (\textit{AMC 10 10B 2008}) A rectangular floor measures $a$ by $b$ feet, where $a$ and $b$ are positive integers with $b > a$. An artist paints a rectangle on the floor with the sides of the rectangle parallel to the sides of the floor. The unpainted part of the floor forms a border of width 1 foot around the painted
rectangle and occupies half of the area of the entire floor. How many possibilities are there for the ordered pair $(a, b)$?


 $\mathrm{(A)}\ 1\qquad\mathrm{(B)}\ 2\qquad\mathrm{(C)}\ 3\qquad\mathrm{(D)}\ 4\qquad\mathrm{(E)}\ 5$


 
 \textbf{Answer: }B



\textbf{Solution}

Because the unpainted part of the floor covers half the area, then the painted rectangle covers half the area as well. Since the border width is 1 foot, the dimensions of the rectangle are $a-2$ by $b-2$. With this information we can make the equation:

 
 \begin{eqnarray*} ab &=& 2\left((a-2)(b-2)\right) \\ ab &=& 2ab - 4a - 4b + 8 \\ ab - 4a - 4b + 8 &=& 0  \end{eqnarray*}Applying \href{https://artofproblemsolving.com/wiki/index.php?title=Simon\%27s\_Favorite\_Factoring\_Trick}{Simon's Favorite Factoring Trick}, we get

 
 \begin{eqnarray*}ab - 4a - 4b + 16 &=& 8 \\ (a-4)(b-4) &=& 8 \end{eqnarray*}

 Since $b > a$, then we have the possibilities $(a-4) = 1$ and $(b-4) = 8$, or $(a-4) = 2$ and $(b-4) = 4$. This allows for 2 possibilities: $(5,12)$ or $(6,8)$ which gives us $\boxed{\textbf{(B)} \: 2}$

 


	\item (\textit{AMC 10 10B 2008}) Quadrilateral $ABCD$ has $AB = BC = CD$, angle $ABC = 70^\circ$ and angle $BCD = 170^\circ$. What is the measure of angle $BAD$?


 $\mathrm{(A)}\ 75\qquad\mathrm{(B)}\ 80\qquad\mathrm{(C)}\ 85\qquad\mathrm{(D)}\ 90\qquad\mathrm{(E)}\ 95$


 
 \textbf{Answer: }C



\textbf{Solution 1 (Cyclic Quadrilateral)}


                \begin{itemize}
                    \item Note: This solution requires the use of cyclic quadrilateral properties but could be a bit time-consuming during the contest.
                \end{itemize}


                To start off, draw a diagram like in solution two and label the points. Create lines $\overline{AC}$ and $\overline{BD}$. We can call their intersection point $Y$. Note that triangle $BCD$ is an isosceles triangle so angles $CDB$ and $CBD$ are each $5$ degrees. Since $AB$ equals $BC$, angle $BAC$ equals $55$ degrees, thus making angle $AYB$ equal to $60$ degrees. We can also find out that angle $CYB$ equals $120$ degrees. 

 Extend $\overline{CD}$ and $\overline{AB}$ and let their intersection be $E$. Since angle $BEC$ plus angle $CYB$ equals $180$ degrees, quadrilateral $YCEB$ is a cyclic quadrilateral. 

 Next, draw a line from point $Y$ to point $E$. Since angle $YBC$ and angle $YEC$ point to the same arc, angle $YEC$ is equal to $5$ degrees. Since $EYD$ is an isosceles triangle (based on angle properties) and $YAE$ is also an isosceles triangle, we can find that $YAD$ is also an isosceles triangle. Thus, each of the other angles is $\frac{180-120}{2}=30$ degrees. Finally, we have angle $BAD$ equals $30+55=\boxed{85}$ degrees.

 



\textbf{Solution 2}

First, connect the diagonal $DB$, then, draw line $DE$ such that it is congruent to $DC$ and is parallel to $AB$. Because triangle $DCB$ is isosceles and angle $DCB$ is $170^\circ$, the angles $CDB$ and $CBD$ are both $\frac{180-170}{2} = 5^\circ$. Because angle $ABC$ is $70^\circ$, we get angle $ABD$ is $65^\circ$. Next, noticing parallel lines $AB$ and $DE$ and transversal $DB$, we see that angle $BDE$ is also $65^\circ$, and subtracting off angle $CDB$ gives that angle $EDC$ is $60^\circ$.

 Now, because we drew $ED = DC$, triangle $DEC$ is equilateral. We can also conclude that $EC=DC=CB$ meaning that triangle $ECB$ is isosceles, and angles $CBE$ and $CEB$ are equal.

 Finally, we can set up our equation. Denote angle $BAD$ as $x^\circ$. Then, because $ABED$ is a parallelogram, the angle $DEB$ is also $x^\circ$. Then, $CEB$ is $(x-60)^\circ$. Again because $ABED$ is a parallelogram, angle $ABE$ is $(180-x)^\circ$. Subtracting angle $ABC$ gives that angle $CBE$ equals $(110-x)^\circ$. Because angle $CBE$ equals angle $CEB$, we get 
 \[x-60=110-x\], solving into $x=\boxed{85^\circ}$.

 

                    \begin{figure}[H]
                        \centering
                        \includegraphics[width=0.5\textwidth]{Solution/AMC_10_10B_2008_24_Solution_Figure_1.png}
                    \end{figure}
                    

 Side note: this solution was inspired by some basic angle chasing and finding some 60 degree angles, which made me want to create equilateral triangles.

 



\textbf{Solution 3(Using Trig.)}


                    \begin{figure}[H]
                        \centering
                        \includegraphics[width=0.5\textwidth]{Solution/AMC_10_10B_2008_24_Solution_Figure_2.png}
                    \end{figure}
                    

 Let the unknown $\angle BAD$ be $x$.

 First, we draw diagonal $BD$ and $AC$.$I$ is the intersection of the two diagonals. The diagonals each form two isosceles triangles, $\triangle BCD$ and $\triangle ABC$. 

 Using this, we find: $\angle DBC = \angle CDB = 5^\circ$ and $\angle BAC = \angle BCA = 55^\circ$. Expanding on this, we can fill in a couple more angles. $\angle ABD = 70^\circ - 5^\circ = 65^\circ$, $\angle ACD = 170^\circ - 55^\circ = 115^\circ$, $\angle BIA = \angle CID = 180^\circ - (65^\circ + 55^\circ) = 60^\circ$, $\angle BIC =   \angle AID = 180^\circ - 60^\circ = 120^\circ$.

 We can rewrite $\angle CAD$ and $\angle BDA$ in terms of $x$. $\angle CAD = x - 55^\circ$ and $\angle BDA = 180^\circ - (120^\circ + x - 55^\circ) = 115^\circ - x$.

 Let us relabel $AB = BC = CD = a$ and $AD = b$.

 By Rule of Sines on $\triangle ACD$ and $\triangle ABD$ respectively, $\frac{\sin(\angle CAD)}{a} = \frac{\sin(\angle ACD)}{b}$, and $\frac{\sin(\angle ABD)}{b} = \frac{\sin(\angle BDA)}{a}$

 In a more convenient form, $\frac{\sin(x-55^\circ)}{a} = \frac{\sin(115^\circ)}{b} \implies \frac{a}{b} = \frac{\sin(x-55^\circ)}{\sin(115^\circ)}$ 

 and $\frac{\sin(65^\circ)}{b} = \frac{\sin(115^\circ-x)}{a} \implies \frac{a}{b} = \frac{\sin(115^\circ-x)}{\sin(65^\circ)}$

 $\implies \frac{\sin(115^\circ-x)}{\sin(65^\circ)} = \frac{\sin(x-55^\circ)}{\sin(115^\circ)}$

 Now, by identity $\sin(\theta) = \sin(180^\circ-\theta)$, $\sin(65^\circ) = \sin(115^\circ)$

 Therefore, $\sin(115^\circ-x) = \sin(x-55^\circ).$ This equation is only satisfied by option $\boxed{\text {(C) } 85^\circ}$

 Note: I'm pretty bad at Asymptote, if anyone could edit this and fill in the angles into the diagram, that would be pretty cool.

 



\textbf{Solution 4 (Cheese)}

Using a protractor ruler, draw an accurate diagram (\href{https://i.imgur.com/YIafDgs.jpg/img}{Example Diagram}). $\angle BAD$ looks slightly less than $90$ degrees. Therefore the answer is $\boxed{\textbf{(C) } 85}$ as $85$ is slightly less than $90$.

 


 added a single R to spell “ruler” correctly ~ethanhansummerfun

 added a space between the "less" and "than" "lessthan" ~Evanjzhang

 Note: protractors are not allowed during this test so you are basically cheating but who cares ~BOTNATE`

 



\textbf{Solution 5 (annoying amounts of algebra + trig identities)}

place A at the origin of a coordinate system, with D on the x-axislet angle BAD be $\theta$, and AB=BC=CD=1

 The y value of from B-A is $\sin\theta$. The y value from C-B is $\theta-(180^\circ-70^\circ)=\theta-110^\circ$. The y value from D-C is $\theta-110^\circ-(180^\circ-170^\circ)=\theta-120^\circ$

 The angles for the vectors from B to C and C to D are angle\_original-(180-angle\_polygon) are because the external angle of the polygon is 180-external angle, which is subtracted from the angle since it heads that amount off from the original direction. 

 since D-C+C-B+B-A=D-A=0 (since A, D are both on x-axis and have the same y value of 0), then: 
 \[\sin\theta+\sin(\theta-110^\circ)+\sin(\theta-120^\circ)=0\]

 from here we expand out the trig expressions using sin addition and isolate $\theta$

 
 \[\sin\theta+\sin(\theta-110^\circ)+\sin(\theta-120^\circ)=0\]
 \[\sin\theta+\sin(\theta)\cos(110^\circ)-\cos(\theta)\sin(110^\circ)   +   \sin(\theta)\cos(120^\circ)-\cos(\theta)\sin(120^\circ)=0\]

 


 
 \[1+\cos(110^\circ)-\cot(\theta)\sin(110^\circ)+\cos(120^\circ)-\cot(\theta)\sin(120^\circ)=0\]
 \[\cot(\theta)=\frac{1+\cos(110^\circ)+\cos(120^\circ)}{\sin(110^\circ)+\sin(120^\circ)}\]

 At this point if you are a human calculator feel free to solve, otherwise we want to try and evaluate the right hand side into some nice expression (ideally cot of an angle).

 
 \[\cot(\theta)=\frac{1+\cos(120^\circ-10^\circ)+\cos(120^\circ)}{\sin(120^\circ-10^\circ)+\sin(120^\circ}\]

 
 \[\cot(\theta)=\frac{1+\cos(120^\circ)\cos(10^\circ)+\sin(120^\circ)\sin(10^\circ)+\cos(120^\circ)} {\sin(120^\circ)\cos(10^\circ)-\cos(120^\circ)\sin(10^\circ)+\sin(120^\circ}\]

 
 \[\cot(\theta)=\frac{1+\frac{-1}{2}\cos(10^\circ)+\frac{\sqrt{3}}{2}\sin(10^\circ)+\frac{-1}{2}} {\frac{\sqrt{3}}{2}\cos(10^\circ)-\frac{-1}{2}\sin(10^\circ)+\frac{\sqrt{3}}{2}}\]

 
 \[\cot(\theta)=\frac{1-\cos(10^\circ)+\sqrt{3}\sin(10^\circ)} {\sqrt{3}\cos(10^\circ)+\sin(10^\circ)+\sqrt{3}}\]

 since the expression still isn't simplified, notice that using the double angle identity on cosine can be used to cancel the 1, and $\sqrt{3}$

 Let: $\cos(5^\circ)=a$, $\sin(5^\circ)=b$

 
 \[\cot(\theta)=\frac{1-(1-2b^2)+\sqrt{3}\cdot(2ab)} {\sqrt{3}(2a^2-1)+2ab+\sqrt{3}}\]

 
 \[\cot(\theta)=\frac{2b^2+2ab\sqrt{3}} {2a^2\sqrt{3}+2ab}\]

 
 \[\cot(\theta)=\frac{2b(b+a\sqrt{3})} {2a(a\sqrt{3}+b)}\]

 
 \[\cot(\theta)=\frac{b}{a}\]

 
 \[\cot(\theta)=\tan(5^\circ)\]

 
 \[\cot(\theta)=\cot(85^\circ)\]

 



\textbf{Solution 6 (guess and check)}

Obtain: $\sin\theta+\sin(\theta-110^\circ)+\sin(\theta-120^\circ)=0$ from solution 5We now guess $\theta=85^\circ$ and try to verify

 
 \[\sin85^\circ+\sin-25^\circ+\sin-35^\circ=0\]
 \[\cos5^\circ=\sin25^\circ+\sin35^\circ\]
 \[\cos5^\circ=\sin(30^\circ-5^\circ)+\sin(30^\circ+5^\circ)\]
 \[\cos5^\circ=2\sin30^\circ\cos5^\circ\]
 \[\cos5^\circ=\cos5^\circ\]

 



\textbf{Solution 7 (alternate way to bash out algebra + trig identities)}

Use equation from Solution 5: 
 \[\sin\theta+\sin(\theta-110^\circ)+\sin(\theta-120^\circ)=0\]
 \[\sin\theta+\sin(\theta-115^\circ+5^\circ)+\sin(\theta-115^\circ-5^\circ)=0\]
 \[\sin\theta+2\sin(\theta-115^\circ)\cos(5^\circ)=0\]
 \[\sin\theta-2\sin(\theta+65^\circ)\cos(5^\circ)=0\]
 \[\sin\theta=2\sin(\theta+65^\circ)\cos(5^\circ)\]

 now guess $\theta=85$ so that the cos(5) is dealt with (and then check it works)

 If you refuse to guess work through the following algebra :D 

 Let: $\phi=\theta+65^\circ$

 
 \[\sin(\phi-65^\circ)=2\sin\phi\cos(5^\circ)\]
 \[\sin\phi\cos(65^\circ)-\cos\phi\sin(65^\circ)=2\sin\phi\cos(5^\circ)\]
 \[\cos(65^\circ)-\cot(\phi)\sin(65^\circ)=2\cos(5^\circ)\]
 \[\cot(\phi)(\sin(60^\circ)\cos(5^\circ)+\cos(60^\circ)\sin(5^\circ))=(\cos(60^\circ)\cos(5^\circ)-\sin(60^\circ)\sin(5^\circ)) -2\cos(5^\circ)\]
 \[\cot(\phi)(\frac{\sqrt{3}}{2}\cos(5^\circ)+\frac{1}{2}\sin(5^\circ))=(\frac{1}{2}\cos(5^\circ)-\frac{\sqrt{3}}{2}\sin(5^\circ)) -2\cos(5^\circ)\]
 \[\cot(\phi)(\frac{\sqrt{3}}{2}\cos(5^\circ)+\frac{1}{2}\sin(5^\circ))=(\frac{-3}{2}\cos(5^\circ)-\frac{\sqrt{3}}{2}\sin(5^\circ))\]
 \[\cot(\phi)(\frac{\sqrt{3}}{2}\cos(5^\circ)+\frac{1}{2}\sin(5^\circ))=-\sqrt{3} (\frac{\sqrt{3}}{2}\cos(5^\circ)+\frac{1}{2}\sin(5^\circ))\]
 \[\cot(\phi)=-\sqrt3\]

 
 \[\phi=150^\circ\]

 

 \[\theta=150^\circ-65^\circ=85^\circ\]

 
 \[\angle BAD = 85^\circ\]

 


	\item (\textit{AMC 10 10B 2008}) Michael walks at the rate of $5$ feet per second on a long straight path. Trash pails are located every $200$ feet along the path. A garbage truck travels at $10$ feet per second in the same direction as Michael and stops for $30$ seconds at each pail. As Michael passes a pail, he notices the truck ahead of him just leaving the next pail. How many times will Michael and the truck meet? 


 $\mathrm{(A)}\ 4\qquad\mathrm{(B)}\ 5\qquad\mathrm{(C)}\ 6\qquad\mathrm{(D)}\ 7\qquad\mathrm{(E)}\ 8$


 
 \textbf{Answer: }B



\textbf{Solution 1}

Pick a coordinate system where Michael's starting pail is $0$ and the one where the truck starts is $200$.Let $M(t)$ and $T(t)$ be the coordinates of Michael and the truck after $t$ seconds.Let $D(t)=T(t)-M(t)$ be their (signed) distance after $t$ seconds.Meetings occur whenever $D(t)=0$.We have $D(0)=200$. 

 The truck always moves for $20$ seconds, then stands still for $30$. During the first $20$ seconds of the cycle the truck moves by $200$ feet and Michael by $100$, hence during the first $20$ seconds of the cycle $D(t)$ increases by $100$.During the remaining $30$ seconds $D(t)$ decreases by $150$.

 From this observation it is obvious that after four full cycles, i.e. at $t=200$, we will have $D(t)=0$ for the first time.

 During the fifth cycle, $D(t)$ will first grow from $0$ to $100$, then fall from $100$ to $-50$. Hence Michael overtakes the truck while it is standing at the pail. 

 During the sixth cycle, $D(t)$ will first grow from $-50$ to $50$, then fall from $50$ to $-100$. Hence the truck starts moving, overtakes Michael on their way to the next pail, and then Michael overtakes the truck while it is standing at the pail. 

 During the seventh cycle, $D(t)$ will first grow from $-100$ to $0$, then fall from $0$ to $-150$. Hence the truck meets Michael at the moment when it arrives to the next pail.

 Obviously, from this point on $D(t)$ will always be negative, meaning that Michael is already too far ahead. Hence we found all $\boxed{5 \Longrightarrow B}$ meetings.

 The movement of Michael and the truck is plotted below: Michael in blue, the truck in red.

 
                    \begin{figure}[H]
                        \centering
                        \includegraphics[width=0.5\textwidth]{Solution/AMC_10_10B_2008_25_Solution_Figure_1.png}
                    \end{figure}
                    

 



\textbf{Solution 2}

The truck takes $20$ seconds to go from one pail to the next and then stops for $30$ seconds at the new pail. Thus it sets off from a pail every 50 sec. Let $t$ denote the time elapsed and write $t=50k + \Delta$, where $\Delta \in [0,50)$. In  this time Michael has traveled $5t = 250k+5\Delta$ feet. What about the truck? In the first $50k$ seconds the truck covers $k$ pails, i.e. $200k$ feet so it moves $200+200k$ feet from Michael's starting point. Then we have two cases: 

 \textbf{(a)} if $\Delta < 20$, then the truck travels an additional $10\Delta$ feet. For them to intersect we must have $200+200k + 10\Delta = 250k + 5\Delta$. Solving, we get $\Delta = 10k - 40$. Since $\Delta$ must lie in the interval $[0,20)$ we get $k \in \{4, 5\}$.

 \textbf{(b)} if $\Delta \in [20, 50)$, then the truck travels an additional $200$ feet. For them to intersect we must have $200+200k + 200 = 250k + 5\Delta$. Solving, we get $\Delta = 10(8-k)$. Since $\Delta$ must lie in the interval $[20,50)$ we get $k \in \{4, 5, 6\}$. 

 Thus Michael intersects with the truck $5$ times, which is option $\boxed{\textbf{(B)}}$.

 



\textbf{Solution 3}

We make a chart by seconds in increments of ten. 

 $\begin{tabular}{c|c|c}Time (s) &Distance of Michael&Distance of Garbage\\ \hline 0&0&200\\ 10&50&300\\ 20&100&400\\ 30&150&400\\ 40&200&400\\ 50&250&400\\ 60&300&500\\ 70&350&600\\ 80&400&600\\ 90&450&600\\ 100&500&600\\ 110&550&700\\ 120&600&800\\ 130&650&800\\ 140&700&800\\ 150&750&800\\ 160&800&900\\ 170&850&1000\\ 180&900&1000\\ 190&950&1000\\ 200&1000&1000\\ 210&1050&1100\\ 220&1100&1200\\ 230&1150&1200\\ 240&1200&1200\\ 250&1250&1200\\ 260&1300&1300\\ 270&1350&1400\\ 280&1400&1400\\ 290&1450&1400\\ 300&1500&1400\\ 310&1550&1500\\ 320&1600&1600\\ 330&1650&1600\\ 340&1700&1600\\ 350&1750&1600\\ 360&1800&1700\\ 370&1850&1800\\ 380&1900&1800\\ 390&1950&1800\\ 400&2000&1800\\ \end{tabular}$

 
Notice at 200, 240, 260, 280, and 320 seconds, Michael and the garbage truck meet. It is clear that they met at these times, and will meet no more. Thus the answer is $\boxed{\textbf{(B)}}$.

 


 



\textbf{Solution 4 (Educated Guess)}

We can record the time of the truck and the time of Michael at each trash pail. 

 $\begin{tabular}{c|c|c}Distance (s) &Truck Time&Michael Time\\ \hline 400&20, 50&40\\ 600&70, 100&80\\ 800&120, 150&120\\ 1000&170, 200&160\\ 1200&220, 250&200\\ 1400&270, 300&240\\ 1600&320, 350&320\\ \end{tabular}$

 Since Micheal moves faster than the truck, the truck will not catch up to Michael afterwards.

 We see that there are 4 correct times and we add 1 because they also pass between pails once which is option $\boxed{\textbf{(B)}}$

 


\end{enumerate}